
\section{Aplicacions i cardinals}
\iniciTest{t2}{b,d,c,b,a,c,c,a,d,b}

\begin{enumerate}
\pregunta Quina de les seg\"{u}ents relacions $R$ entre $A$ i $B$ \'{e}s una
aplicaci\'{o} d'$A$ en $B$?

\begin{enumerate}
\opcio{a}{$A=\left\{ 1,2,3,4,5\right\} $, $B=\left\{ 1,2\right\} $ i $R=\left\{
(1,1),(3,2),(5,1),(4,1)\right\} $
}
\opcio{b}{$A=B=\mathbb{R}$ i $R=\left\{ (x,y)\in \mathbb{R}\times \mathbb{R}%
:x+y^{3}=0\right\} $
}
\opcio{c}{$A=B=\mathbb{R}$ i $R=\left\{ (x,y)\in\mathbb{R}\times\mathbb{R}%
:xy=1\right\} $
}
\opcio{d}{$A=B=\mathbb{R}$ i $R=\left\{ (x,y)\in\mathbb{R}\times\mathbb{R}:y=%
\sqrt{x-1}\right\} $
}
\end{enumerate}
\feedback

\pregunta Es defineix l'aplicaci\'{o} $f:\mathbb{R}~\longrightarrow~\mathbb{R}$
mitjan\c{c}ant $f(x)=x^{2}+4$. Llavors,

\begin{enumerate}
\opcio{a}{$f^{-1}(\left\{ 0\right\} )=\{-2,2\}$
}
\opcio{b}{$f\left( [0,1]\right) =[0,1]$
}
\opcio{c}{$f^{-1}\left( (0,4)\right) =(0,2)$
}
\opcio{d}{Cap de les anteriors \'{e}s certa
}
\end{enumerate}
\feedback

\pregunta La gr\`{a}fica d'una aplicaci\'{o} $f:\mathbb{R}~\longrightarrow~%
\mathbb{R}$ \'{e}s\FRAME{dtbpF}{2.9447in}{2.0072in}{0pt}{}{}{set11.jpg}{}llavors:

\begin{enumerate}
\opcio{a}{$f$ \'{e}s injectiva
}
\opcio{b}{$f$ \'{e}s exhaustiva
}
\opcio{c}{$f:[0,+\infty )~\longrightarrow ~[0,1]$ \'{e}s bijectiva
}
\opcio{d}{$f:[0,+\infty)~\longrightarrow~\mathbb{R}$ \'{e}s exhaustiva
}
\end{enumerate}
\feedback

\pregunta Considerem les aplicacions: $f:\mathbb{R}-\{-2\}~\longrightarrow ~%
\mathbb{R}-\{3\}$ definida per%
\begin{equation*}
f(x)=\frac{3+3x}{x+2}
\end{equation*}
i $g:\mathbb{R}-\{3\}~\longrightarrow~\mathbb{R}$ definida per $g(x)=x^{3}$.
Llavors, quina de les seg\"{u}ents afirmacions \'{e}s falsa?

\begin{enumerate}
\opcio{a}{$f$ \'{e}s bijectiva i $f^{-1}(x)=\frac{3-2x}{x-3}$
}
\opcio{b}{$g$ \'{e}s bijectiva i $g^{-1}(x)=\sqrt[3]{x}$
}
\opcio{c}{$f\circ g$ no \'{e}s aplicaci\'{o}
}
\opcio{d}{$g\circ f$ \'{e}s injectiva i $(g\circ f)(x)=\left( \frac{3+3x}{x+2}%
\right) ^{3}$
}
\end{enumerate}
\feedback

\pregunta Sigui $f:A~\longrightarrow~B$ una aplicaci\'{o} i suposem que $\#A=n$
i $\#B=m$. Llavors:

\begin{enumerate}
\opcio{a}{Si $f$ \'{e}s injectiva, llavors $n\leq m$
}
\opcio{b}{Si $f$ \'{e}s exhaustiva, llavors $n\leq m$
}
\opcio{c}{Si $n=m$, llavors $f$ \'{e}s bijectiva
}
\opcio{d}{No pot oc\'{o}rrer que $n<m$
}
\end{enumerate}
\feedback

\pregunta Si $f,g$ s\'{o}n aplicacions de $\mathbb{R}$ en $\mathbb{R}$ tals que $%
g(x)=x^{3}$ i $(g\circ f)(x)=x^{3}-3x^{2}+3x-1$. Llavors:

\begin{enumerate}
\opcio{a}{$f(x)=x+1$
}
\opcio{b}{$f(x)=1-x$
}
\opcio{c}{$f(x)=x-1$
}
\opcio{d}{$f(x)=-1-x$
}
\end{enumerate}
\feedback

\pregunta Sigui $f:A~\longrightarrow~B$ una aplicaci\'{o} i considerem $%
X,Y\subset A$ i $Z,T\subset B$, llavors

\begin{enumerate}
\opcio{a}{$f(X\cap Y)=f(X)\cap f(Y)$
}
\opcio{b}{$f^{-1}\left( f(X)\right) =X$
}
\opcio{c}{$f^{-1}(Z\cap T)=f^{-1}(Z)\cap f^{-1}(T)$
}
\opcio{d}{$f\left( f^{-1}(Z)\right) =Z$
}
\end{enumerate}
\feedback

\pregunta Efectuant una mostra de 1000 individus s'observa que mengen peix i
carn per\`{o} no ous 60, peix i ous per\`{o} no carn 40, carn i ous per\`{o}
no peix 30, nom\'{e}s peix 50, nom\'{e}s carn 40 i nom\'{e}s ous 30.
Tots mengen carn, ous o peix. Quants mengen peix?

\begin{enumerate}
\opcio{a}{900
}
\opcio{b}{750
}
\opcio{c}{800
}
\opcio{d}{Cap de les anteriors
}
\end{enumerate}
\feedback

\pregunta En una classe de 100 alumnes que s'han examinat de matem\`{a}tiques i F%
\'{\i}sica es coneixen els seg\"{u}ents resultats: No han aprovat cap
assignatura 20 alumnes. Han aprovat les dues assignatures 25 alumnes. Han
aprovat el doble d'alumnes Matem\`{a}tiques que F\'{\i}sica.
Quants alumnes han aprovat Matem\`{a}tiques?

\begin{enumerate}
\opcio{a}{10
}
\opcio{b}{20
}
\opcio{c}{35
}
\opcio{d}{70
}
\end{enumerate}
\feedback

\pregunta En el conjunt dels nombres naturals menors que 500, %
quants nombres no són m\'{u}ltiples de 2, ni de 3, ni de 5?

\begin{enumerate}
\opcio{a}{120
}
\opcio{b}{134
}
\opcio{c}{100
}
\opcio{d}{Cap de les anteriors
}
\end{enumerate}
\feedback
\end{enumerate}

